\section{Asymptotic formulas on Magnitude- and Hessian- pruning}\label{sec prune risk}
Here, we use Theorem \ref{thm:master_W2} to characterize the risk of the magnitude- and Hessian- pruned solutions. This section supplements the discussion in Section \ref{sec:risk}. For completeness, first, we recall the magnitude-pruning results of Section \ref{sec:risk} and restate Corollary \ref{cor:mag} below. This corollary characterizes the performance of magnitude pruning. Following this, we shall further discuss Hessian pruning.


\subsection{Magnitude-based pruning}

We begin with the following necessary definitions.  Define the hard-thresholding function with fixed threshold $t\in\R_+$ as follows:
\begin{align}\label{eq:threshold_app}
\Tc_t(x) = \begin{cases} x & \text{if } |x|>t \\ 0 & \text{otherwise} \end{cases}.
\end{align}
Further, {given model sparsity target $1>\alpha>0$}, define the threshold $t^\star$ as follows:
\begin{align}\label{eq:tstar_app}
t^\star:=\sup\left\{t\in\R\,:\, \Pr(|X_{\kappa,\sigma^2}|\geq t) \geq \alpha \right\}.
\end{align}


\cormag*

The proof of the corollary above, is given in Section \ref{sec:risk}. Below, we extend the results to Hessian-based pruning.


\subsection{Hessian-based pruning}
Let $\bth$ be the min-norm solution in \eqref{eq:min_norm}. Recall that the Hessian-pruned model (via Optimal Brain Damage) $\betab_s^H$ at sparsity $s$ is given by
\begin{align}\label{eq:hp}
\betab_s^H = \hat\Sigmab^{-1/2}\Tb_s(\hat\Sigmab^{1/2}\betab),
\end{align}
where $\hat\Sigmab=\diag{\X^T\X}/n$ the diagonal of the empirical covariance matrix. 

We will argue that the following formula characterizes the asymptotic risk of the Hessian pruning solution. 
Recall the notation in \eqref{eq:threshold_app} and define
\begin{align}\label{eq:tdiam_app}
t^\diamond:=\sup\left\{t\in\R\,:\, \Pr(|\Lambda^{1/2}X_{\kappa,\sigma^2}|\geq t) \geq \alpha \right\}.
\end{align}
The risk of the Hessian-pruned model satisfies the following in the limit of $n,p,s\rightarrow\infty$ at rates $\kappa:=p/n>1$ and $\alpha:=s/p\in(0,1)$ (cf. Assumption \ref{ass:linear}):
\begin{align}\label{eq:hessian_risk}
\Lc(\bth^H_s) \rP \sigma^2 + \E\left[ \left(\Lambda^{1/2}B-\Tc_{t^\diamond}(\Lambda^{1/2}X_{\kappa,\sigma^2})\right)^2 \right],
\end{align}
where the expectation is over $(\Lambda,B,H)\sim\mu\otimes\Nn(0,1)$. In our LGP experiments, we used this formula \eqref{eq:hessian_risk} to accurately predict the Hessian-based pruning performance. 
%\begin{corollary}[Risk of Hessian-pruning]\label{cor:hes}
%Let the same assumptions and notation as in the statement of Theorem \ref{thm:master_W2} hold. Specifically, let $\hat\betab$ be the min-norm solution in \eqref{eq:min_norm} and $\hat\beta_s^H$ the Hessian-pruned model at sparsity $s$. Recall the notation in \eqref{eq:threshold_app} and define
%\begin{align}\label{eq:tdiam_app}
%t^\diamond:=\inf\left\{t\in\R\,:\, \Pr(|\Lambda^{1/2}X_{\kappa,\sigma^2}|\geq t) \geq \alpha \right\}.
%\end{align}
%The risk of the magnitude-pruned model satisfies the following in the limit of $n,p,s\rightarrow\infty$ at rates $\kappa:=p/n>1$ and $\alpha:=s/p\in(0,1)$ (cf. Assumption \ref{ass:linear}):
%\begin{align}
%\Lc(\bth^H_s) \rP \sigma^2 + \E\left[ \Lambda\left(B-\Lambda^{-1/2}\Tc_{t^\diamond}(\Lambda^{1/2}X_{\kappa,\sigma^2})\right) \right],\nn
%\end{align}
%where the expectation is over $(\Lambda,B,H)\sim\mu\otimes\Nn(0,1).$
%\end{corollary}
%\begin{proof}

Recall the definition of the hard-thresholding operator $\Tc_t(x)$. Similar to Section \ref{sec:risk}, we consider a threshold-based pruning vector $$\bth^{\Tc,H}_{t} := \hat\Sigmab^{-1/2}\Tc_{t/\sqrt{p}}(\hat\Sigmab^{1/2}\hat\bt),$$ where $\Tc_t$ acts component-wise. Further define 
$$\bth^{\Tc,H^\star}_{t} := \Sigmab^{-1/2}\Tc_{t/\sqrt{p}}(\Sigmab^{1/2}\hat\bt).$$
Note that $\bth^{\Tc,H^\star}_{t}$ uses the true (diagonal) covariance matrix $\Sigmab$ instead of its sample estimate $\hat\Sigmab$. For later reference, note here that $\hat\Sigmab$ concentrates (entry-wise) to $\Sigmab$. Using boundedness of $\Sigmab$ and standard concentration of sub-exponential random variables.


First, we compute the limiting risk of $\bth^{\Tc,H^\star}_{t}$. Then, we will use the fact that $\hat\Sigmab$ concentrates (entry-wise) to $\Sigmab$, to show that the risks of $\bth^{\Tc,H^\star}_{t}$ and $\bth^{\Tc,H}_{t}$ are arbitrarily close as $p\rightarrow\infty$.

%We start by computing the risk of $\bth^{\Tc,H^\star}_{t}$. 
Similar to \eqref{eq:loss_w2}, 
\begin{align}
\Lc(\bth^{\Tc,H^\star}_t) &= \sigma^2 + (\bt^\star-\bth^{\Tc,H^\star}_t)^T\Sigmab(\bt^\star-\bth^{\Tc,H^\star}_t) \nn \\
  &= \sigma^2 + \frac{1}{p}\sum_{i=1}^p\Sigmab_{i,i}\big(\sqrt{p}\betas_i-\lai^{-1/2}\Tc_{t}(\lai^{1/2}\sqrt{p}\hat\betab_i)\big)^2 \nn
    = \sigma^2 + \frac{1}{p}\sum_{i=1}^p\big(\lai^{1/2}\sqrt{p}\betas_i-\Tc_{t}(\lai^{1/2}\sqrt{p}\hat\betab_i)\big)^2 
      \\
    &= \sigma^2 + \frac{1}{p}\sum_{i=1}^p\big(\lai^{1/2}\sqrt{p}\betas_i-\Tc_{t}(\wh_i + \lai^{1/2}\sqrt{p}\betas_i\big)^2 \nn \\
        &=: \sigma^2 + \frac{1}{p}\sum_{i=1}^p\big(\sqrt{p}\blat_i-\Tc_{t}(\wh_i + \blat_i)\big)^2.\label{eq:hesa}
  %\nn\\
%  &\rP \sigma^2 + \E\left[ \Lambda\left(B-\Tc_t(X_{\kappa,\sigma^2})\right) \ri
\end{align}
In the second line above, we used that $\sqrt{p}\Tc_{t/\sqrt{p}}(x)=\Tc_{t}(\sqrt{p}x)$. In the third line, we recalled the change of variable in \eqref{eq:w}, i.e., $\wh$ is the solution to \eqref{eq:PO}. Finally, in the last line we defined $\blat:=\sqrt{\Sigmab}\betas$ (note that this is related to the saliency score $\bla$ defined in \eqref{saliency eq}).

To evaluate the limit of the empirical average in \eqref{eq:hesa}, we proceed as follows. First, we claim that the empirical distribution of $\sqrt{p}\blat$ converges weakly to the distribution of the random variable $B\sqrt{\Lambda}$, where $(\Lambda,B)\sim\mu$. \so{Note that this convergence is already implied by the proof of Theorem \ref{thm:master_W2} by setting the $g$ function to be zero in \eqref{eq:fdef}.}
%\somm{This part is repetitive, can be simplified.}
For an explicit proof, take any bounded Lipschitz test function $\psi:\R\rightarrow\R$ and call $\psi'(x,y):=\psi({\sqrt{x}}y)$. Then, 
\begin{align}
|\psi'(\lai,\sqrt{p}\betas_i)-\psi'(\lai',\sqrt{p}{\betas_i}')| &= |\psi(\sqrt{\lai}\sqrt{p}\betas_i)-\psi(\sqrt{\lai'}\sqrt{p}{\betas_i}')| \leq C |\sqrt{\lai}\sqrt{p}\betas_i - \sqrt{\lai'}\sqrt{p}{\betas_i}'| \nn
\\
&\leq
C |\sqrt{p}\betas_i - \sqrt{p}{\betas_i}'| + C'|\sqrt{p}{\betas_i}'| |\sqrt{\lai}- \sqrt{\lai'}|\nn  \\
&\leq
C (|\sqrt{p}\betas_i - \sqrt{p}{\betas_i}'| + |\sqrt{p}{\betas_i}'| |{\lai}- {\lai'}|)\nn  \\
&\leq
C (1 + \|[\sqrt{p}\betas_i,\lai]\|_2 +\|[\sqrt{p}{\betas_i}',\lai']\|_2 ) \sqrt{|\sqrt{p}\betas_i - \sqrt{p}{\betas_i}'|^2   + |{\lai}- {\lai'}|^2}\nn.
\end{align}
Thus, $\psi'$ is $\rm{PL(2)}$. Hence, from Assumption \ref{ass:mu},
\begin{align}
\frac{1}{p}\sum_{i=1}^p \psi(\sqrt{p}\blat_i) = \frac{1}{p}\sum_{i=1}^p \psi'(\lai,\sqrt{p}\betas_i) \rP  \E[\psi'(B,\sqrt{\Lambda})] =  \E[\psi(B\sqrt{\Lambda})]\label{eq:convvvv}
\end{align}
Besides, from Theorem \ref{thm:master_W2} applied for $f_\Lc(x,y,z)=zy^2$ (i.e., set $g$ the zero function in \eqref{eq:fdef}), we have that
$$
\frac{1}{p}\sum_{i=1}^p \sqrt{p}(\blat_i)^2 \rP  \E[B^2\Lambda].
$$
Therefore, convergence in \eqref{eq:convvvv} actually holds for any $\psi\in\rm{PL}(2)$. Next, observe that, $\wo$ of \eqref{eq:w_n} can be written in terms of $\blat$ via \begin{align}\label{eq:w_n2}
\w_n=\w'(\tau_n,u_n)& = -\left(\Iden+\frac{u_n}{\tau_n}\Sigmab\right)^{-1}\left(\Sigmab^{1/2}\betas-u_n\sqrt{\kappa}\Sigmab\hba\right)\\
&=-\left(\Iden+\frac{u_n}{\tau_n}\Sigmab\right)^{-1}\left(\blat-u_n\sqrt{\kappa}\Sigmab\hba\right).
\end{align}
After this observation, the convergence proof can be finalized by using a modified version of Assumption \ref{ass:mu} as follows.%in two ways. First approach is

{
\begin{assumption}[Empirical distribution for saliency] \label{ass:lambda} Set $\blat_i=\sqrt{\lai}\bts_i$. The joint empirical distribution of $\{(\lai,\blat_i)\}_{i\in[p]}$ converges in {Wasserstein-k} distance to a probability distribution $\mu=\mu(\Lambda,S)$ on $\R_{>0}\times\R$ for some \fx{$k\geq 4$}. That is
$
\frac{1}{p}\sum_{i\in[p]}\delta_{(\lai,\sqrt{p}\blat_i)} \stackrel{W_k}{\Longrightarrow} \mu.
$
%weakly to a probability distribution $\mu=\mu(\Lambda, S)$ on $\R_{>0}\times\R$ with bounded $4$th moment and assume that as $p\rightarrow\infty$, the $4$th moment of the empirical distribution converges to the $4$th moment of $\mu$, i.e.,
%$$
%\frac{1}{p}\sum_{i=1}^p\left(\sqrt{(\sqrt{p}\blat_i)^{2}+\lai^2}\right)^4 \rP \E_{(\Lambda,S)\sim\mu}\left[\left(\sqrt{\Lambda^2+S^2}\right)^4\right]<\infty.
%$$
\end{assumption}
}

Under this assumption, it can be verified that, the exact same proof strategy we used for magnitude-based pruning would apply to Hessian-based pruning by replacing $\bts$ with $\blat$. The reason is that the Hessian pruning risk is a $\text{PL}(4)$ function of $\h,\blat,\bSi$ (e.g.~can be shown in a similar fashion to Lemma \ref{lem:hPL}). Observe that Assumption \ref{ass:lambda} is a reasonable assumption given the naturalness of the saliency score. If we only wish to use the earlier Assumption \ref{ass:mu} rather than Assumption \ref{ass:lambda}, one can obtain the equivalent result by modifying \ref{ass:mu} to enforce a slightly higher order bounded moment and convergence condition. Finally, one needs to address the perturbation due to the finite sample estimation of the covariance. Note that, even if the empirical covariance doesn't converge to the population, its diagonal weakly converges to the population (as we assumed the population is diagonal). The (asymptotically vanishing) deviation due to the finite sample affects can be addressed in an identical fashion to the deviation analysis of $\tau_n,u_n$ at \eqref{eq:ivstep1} and \eqref{eq:betav}. While these arguments are reasonably straightforward and our Hessian pruning formula accurately predicts the empirical performance, the fully formal proof of the Hessian-based pruning is rather lengthy to write and does not provide additional insights. 
%Thus the associated formal theoretical statement is deferred to a future version.

\cmt{
if one assumes Assumption \ref{ass:mu} with the 
The overall function $f(\blat_i,$ becomes
\[
h(x,y,z)=f(g(y,z,x),y,z)=y(x-g(z))^2
\]}
%\cts{Using the assumption of the corollary that $\beta$}
%\end{proof}






